\documentclass{beamer}
\usetheme{Madrid}
\usepackage{amsmath}
\usepackage{amssymb}

\title{Current Rule-Based Behavior Planner}
\subtitle{Lane Safety and Lane Selection Logic}
\author{}
\date{}

\begin{document}

\begin{frame}
  \titlepage
\end{frame}

\begin{frame}{Planner Role}
  \begin{itemize}
    \item The rule-based planner runs inside the main loop and selects the behavior:
    \[
      \texttt{lane\_follow}, \quad \texttt{lane\_change\_left}, \quad \texttt{lane\_change\_right}
    \]
    \item Main inputs:
    \begin{itemize}
      \item lane safety scores \(S_{\text{lane}}\)
      \item ego lane id
      \item route-optimal lane id from the global planner
      \item front-obstacle distance in each lane
      \item ego lateral offset and heading error
      \item driving mode: NORMAL or INTERSECTION
    \end{itemize}
    \item Main output: target lane for the blue-dot / local reference generator.
  \end{itemize}
\end{frame}

\begin{frame}{Obstacle Classification in Ego Frame}
  For each obstacle, the planner computes longitudinal position in the ego frame:
  \[
    s = \cos(\psi_e)(x_o - x_e) + \sin(\psi_e)(y_o - y_e)
  \]
  \begin{itemize}
    \item \(s \ge 0\): obstacle is in front of ego
    \item \(s < 0\): obstacle is behind ego
    \item For each lane, only the nearest front and nearest rear obstacle are used.
  \end{itemize}

  \vspace{0.5em}
  Variables:
  \begin{itemize}
    \item \((x_e, y_e, \psi_e)\): ego position and heading
    \item \((x_o, y_o)\): obstacle position
    \item \(s\): obstacle longitudinal distance relative to ego
  \end{itemize}
\end{frame}

\begin{frame}{Lane Safety Score}
  The obstacle-level score is built from distance, TTC, and TTC trend:
  \[
    \phi_d = \frac{\max(d - d_{\text{safe}}, 0)}{\max(d - d_{\text{safe}}, 0) + d_{\text{safe}}}
  \]
  \[
    \phi_{\text{ttc}} = \frac{\max(\text{TTC} - \text{TTC}_{\text{safe}}, 0)}{\max(\text{TTC} - \text{TTC}_{\text{safe}}, 0) + \text{TTC}_{\text{safe}}}
  \]
  \[
    \phi_m = \frac{1}{1 + e^{-k m}}
  \]
  \[
    S_{\text{obs}} = \phi_d \, \phi_{\text{ttc}} \, \phi_m
  \]

  Hard safety constraint:
  \[
    d < d_{\text{safe}} \quad \text{or} \quad \text{TTC} < \text{TTC}_{\text{safe}}
    \Longrightarrow
    S_{\text{obs}} = 0
  \]
\end{frame}

\begin{frame}{TTC and TTC Trend}
  TTC is computed from closing speed:
  \[
    \text{TTC} =
    \begin{cases}
      \dfrac{d}{\max(\varepsilon, \Delta v)}, & \Delta v > 0 \\
      \infty, & \Delta v \le 0
    \end{cases}
  \]

  \vspace{0.5em}
  The TTC trend is the linear-regression slope:
  \[
    m = \frac{\sum_i (t_i - \bar{t})(\text{TTC}_i - \overline{\text{TTC}})}
             {\sum_i (t_i - \bar{t})^2}
  \]

  \begin{itemize}
    \item \(d\): front or rear longitudinal distance
    \item \(\Delta v\): closing speed
    \item \(m > 0\): TTC is growing, so the situation is becoming safer
    \item \(m < 0\): TTC is shrinking, so risk is increasing
  \end{itemize}
\end{frame}

\begin{frame}{Lane Score}
  For each lane, the planner evaluates:
  \[
    S_{\text{front}} = \text{score of nearest front obstacle}
  \]
  \[
    S_{\text{rear}} = \text{score of nearest rear obstacle}
  \]
  \[
    S_{\text{lane}} = \min(S_{\text{front}}, S_{\text{rear}})
  \]

  \begin{itemize}
    \item If a lane has no relevant obstacle, its score defaults to \(1.0\).
    \item The \(\min(\cdot)\) operator makes the lane conservative:
    one dangerous front or rear interaction can make the whole lane unsafe.
  \end{itemize}
\end{frame}

\begin{frame}{Normal-Mode Decision Logic}
  \textbf{Goal:} stay on the global route unless a blockage forces a detour.

  \begin{enumerate}
    \item Choose the route-optimal lane \(L_r\).
    \item If ego is already in \(L_r\):
    \begin{itemize}
      \item keep following \(L_r\) when it is safe
      \item if there is a front obstacle ahead in \(L_r\), move to an adjacent safe lane
    \end{itemize}
    \item If ego is on a temporary detour lane:
    \begin{itemize}
      \item return to \(L_r\) once \(L_r\) becomes safe again
      \item otherwise continue in the temporary lane
    \end{itemize}
  \end{enumerate}

  Lane change is allowed only if target-lane safety is above a threshold:
  \[
    S_{\text{target}} > S_{\text{threshold}}
  \]
\end{frame}

\begin{frame}{Intersection-Mode Decision Logic}
  \textbf{Goal:} let the global route dominate inside intersections.

  \begin{itemize}
    \item The desired lane is the route-optimal lane.
    \item No safety-based lane preference override is used to change the route.
    \item If a lane change is needed, only one lane step is allowed per decision.
    \item A lane change is started only if the immediate target lane is safe enough:
    \[
      S_{\text{target}} > S_{\text{threshold}}
    \]
  \end{itemize}

  This keeps the blue dot aligned with the global route while still enforcing a
  minimum safety check before lateral motion.
\end{frame}

\begin{frame}{Lane-Change State Machine}
  States:
  \[
    \texttt{IDLE}, \quad \texttt{CHANGING\_LEFT}, \quad \texttt{CHANGING\_RIGHT}
  \]

  Once a lane change starts, it is not interrupted until completion.

  Completion conditions:
  \[
    \text{ego lane} = \text{target lane}
  \]
  \[
    |e_y| \le e_{y,\max}
    \quad \text{and} \quad
    |e_\psi| \le e_{\psi,\max}
  \]

  \begin{itemize}
    \item \(e_y\): lateral offset from lane center
    \item \(e_\psi\): heading error relative to lane direction
    \item \(e_{y,\max}\), \(e_{\psi,\max}\): lane-change completion thresholds
  \end{itemize}
\end{frame}

\begin{frame}{Key Parameters and Interpretation}
  Important planner parameters:
  \begin{itemize}
    \item \(d_{\text{safe}}\): minimum acceptable longitudinal gap
    \item \(\text{TTC}_{\text{safe}}\): minimum acceptable time-to-collision
    \item \(k\): gain in the TTC-slope sigmoid
    \item \(S_{\text{threshold}}\): minimum target-lane safety for lane change
    \item \(\Delta_{\text{hyst}}\): hysteresis margin before leaving current lane
  \end{itemize}

  \vspace{0.5em}
  Summary:
  \begin{itemize}
    \item lane safety converts obstacle interactions into a score in \([0,1]\)
    \item the behavior planner is route-first, not safety-score-only
    \item lane changes are conservative, one-step, and threshold-gated
  \end{itemize}
\end{frame}

\end{document}
